% Ad Familiares 15.14 (ad-familiares-15-14) --- English translation
% Translated by Claude (Anthropic) from the Latin (Perseus).
% Style guide: STYLE.md.

\ciceroLetterOpener{M. CICERO IMP. S. D. C. CASSIO PROQ.}

\ciceroSection{1} That you give me M. Fadius as a friend on your recommendation --- I make no profit out of that. For he has long been on my books, and is dear to me on account of his great civility and his attentiveness; but still, because I have perceived how unusually you regard him, I have become much more his friend than I was. And so, although your letter did its work, far more weight of recommendation with me lay in the spirit of the man himself toward you, which I have seen for myself and known.

\ciceroSection{2} But concerning Fadius we shall do zealously what you ask. As for you: I should have wished, for many reasons, that we could have met --- first, that I might see, after so long an interval, the man whom I have long held in the highest regard; next, that I might congratulate you in person, as I have done by letter; further, that we might share with each other, you your concerns, I mine, what each of us wished to communicate; and last, that our friendship --- cultivated on both sides with the fullest good offices, but interrupted in its day-to-day exchange by long stretches of time --- might be reinforced more vigorously.

\ciceroSection{3} Since that has not happened, we shall use the resource of letters and accomplish at a distance about the same things we should accomplish in person. There is one pleasure of the heart, of course, that lies in seeing you, and that cannot be had by letter; another is congratulation, which is indeed a thinner thing on the page than it would be if I were looking at you and offering it; but I have offered it before and offer it now, and I congratulate you both on the magnitude of the things you have done and on the timeliness of the moment, in that on your departure from your province the highest praise and the highest goodwill of the province escorted you.

\ciceroSection{4} The third item is that what we should have communicated to each other about my own affairs in person, we should accomplish by letter all the same. For other reasons I strongly advise you to make speed to Rome. The matters I left were tranquil where you were concerned, and I understand that this recent victory of yours will make your arrival a brilliant one. But if there are burdens on your own people: if they are such that you can bear them, hurry --- nothing will be more elegant for you, nothing more glorious. But if they are greater, consider lest your arrival fall at a wholly unsuitable time. The whole calculation of this is yours; you know what you can bear. If you can, the thing is praiseworthy and popular; if you plainly cannot, in your absence you will the more easily bear men's talk.

\ciceroSection{5} As for me, I press upon you in this letter the same thing I pressed in the last: that you strain every nerve to see that no extension of time be given me for this province, which both the Senate and the people willed to be annual. So earnestly do I urge this upon you that I count my fortunes as resting upon it. You have our Paullus, who is most eager for my interest; there is Curio, and there is Furnius. I should wish you to exert yourself, then, as though everything for me lay in this matter.

\ciceroSection{6} The last item among those I had set out is the confirmation of our friendship. About this no abundance of words is needed. You sought me out when you were a boy; I, for my part, have always reckoned that you would be an adornment to me. You were also a protection in my most grievous times. After your departure there grew up between me and your Brutus a very close familiarity. And so in the talent and industry of you both I think I have a great store both of pleasure and of dignity laid up for me. That you would confirm this by your own zeal, I earnestly ask: send me letters both at once and continuously, and, when you reach Rome, as often as you possibly can.
